\begin{abstract}
The integration of Artificial Intelligence (AI) into the Smart Grid is critical for optimizing energy distribution, forecasting demand, and detecting anomalies in real-time. However, the deployment of these complex models introduces significant security vulnerabilities. While traditional adversarial attacks target model integrity to induce misclassifications, this report investigates ``Sponge Attacks''—a class of denial-of-service (DoS) attacks targeting model availability. Sponge attacks exploit the computational density of deep learning models to exhaust system resources, specifically maximizing inference latency and energy consumption—including through maximizing bit-flip transitions that directly increase dynamic power dissipation. In the context of the Smart Grid, such delays can destabilize grid operations or drain battery-powered edge devices like smart meters. This paper conducts a meta-analysis of existing literature on sponge examples and poisoning, identifies vulnerable components within smart grid architectures, and presents a proof-of-concept experiment demonstrating the impact of these attacks on model performance. Finally, we explore mitigation strategies to secure grid AI systems against resource exhaustion.
\end{abstract}

\begin{IEEEkeywords}
Smart Grid, Sponge Attacks, Adversarial Machine Learning, Energy-Latency, Availability Attacks, Deep Learning Security.
\end{IEEEkeywords}
